% Authored: implemented security controls.

\chapter{Security}
\label{ch:security}

This chapter documents controls that exist. \srcfile{SECURITY.md} states the
project's own security model: the default demo runs locally without API keys or
managed services; optional integrations must be configured explicitly; and real
customer data, secrets, keys, or private deployment configuration must not
appear in issues, pull requests, examples, or tests.

\section{Application controls}

\begin{description}
  \item[API-key authorization] Optional, disabled by default. Three ordered
    roles --- \texttt{reader} (rank 1), \texttt{writer} (2), \texttt{admin}
    (3) --- are checked against the \texttt{X-Feedback-Agent-Key} header. A
    higher-ranked key satisfies a lower-ranked requirement. Read routes require
    \texttt{reader}, mutating workflow routes require \texttt{writer}, and
    indexing and ingestion-job submission require \texttt{admin}; the required
    role for every route appears in \cref{ch:http}. Enabling authorization with
    no key configured raises \texttt{ApiAuthConfigurationError}, which maps to
    \texttt{500} --- a misconfiguration fails loudly rather than silently
    allowing traffic.
  \item[Rate limiting] Optional, disabled by default. An in-process
    fixed-window limiter over all routes except the probes, keyed by API key
    when present and by client address otherwise. It is per worker, not per
    service; see \cref{ch:runtime}.
  \item[CORS] Explicit origin allow-list, defaulting to the local Vite
    development and preview ports. A wildcard origin automatically disables
    credentialed requests, which prevents the unsafe wildcard-plus-credentials
    combination.
  \item[Prompt-injection guardrails] Two independent gates; see
    \cref{adr:008}.
  \item[PII redaction] Deterministic pattern-based redaction of email
    addresses, phone numbers, and access tokens, applied on the data path
    before storage and indexing. The token patterns cover common credential
    prefixes (\texttt{sk}, \texttt{pk}, \texttt{ghp}, \texttt{github\_pat},
    Slack \texttt{xox*}, \texttt{AKIA}/\texttt{ASIA}), assignment-shaped
    secrets (\texttt{api\_key = ...}), and bearer tokens. Redaction replaces
    matches with explicit placeholders rather than deleting them, so the
    redaction is visible in the stored text.
  \item[Input validation] Every request body, CSV row, and stream event is
    validated against a Pydantic model at the boundary; see \cref{adr:004}.
  \item[Error mapping] Domain exceptions are converted to status codes at the
    HTTP boundary, so stack traces and internal messages are not returned to
    callers.
\end{description}

\section{Supply-chain controls}

\begin{description}[nosep]
  \item[Dependency locking] \srcfile{poetry.lock} and
    \srcfile{frontend/package-lock.json} are both version controlled; CI
    installs from them.
  \item[Dependency review] Every pull request is scanned by
    \texttt{actions/dependency-review-action}, failing on high severity.
  \item[npm audit] The frontend workflow runs \texttt{npm audit} as a build
    step.
  \item[Dependabot] Weekly grouped updates for pip, npm, and GitHub Actions.
  \item[CodeQL] Static analysis for Python and JavaScript/TypeScript on
    \texttt{main}, on pull requests, and weekly.
  \item[Least-privilege workflows] Every workflow declares an explicit minimal
    \texttt{permissions} block, and no workflow references any secret.
\end{description}

\section{Vulnerability reporting}

\srcfile{SECURITY.md} directs reporters to GitHub Security Advisories rather
than public issues, and asks for affected component, reproduction, impact, and
known mitigations. Security fixes are stated to apply to \texttt{main} and the
latest release.

\section{Boundaries}

Three limits follow directly from the design and should not be mistaken for
gaps in this document:

\begin{itemize}[nosep]
  \item Authorization is a shared-secret API key, not an identity system. There
    are no users, no sessions, no token expiry, and no revocation beyond
    rotating the configured key.
  \item Guardrails and redaction are deterministic pattern matching. They catch
    the shapes they encode and nothing else.
  \item The rate limiter is in-process, so it bounds a worker rather than a
    deployment.
\end{itemize}
